\documentclass[a4paper]{article}

% \usepackage{german}
% \usepackage{fancyheadings}
\usepackage{a4}
\usepackage{xspace}
\usepackage{longtable}
\usepackage[english]{babel,varioref}
\usepackage{bnf}
\usepackage[dvips]{graphics}
\usepackage[dvips]{epsfig}
\usepackage{times}
% \usepackage{floatflt, epsfig}

\font\fmilogofont=fmilogo15 scaled\magstep0    
\font\skriptuni=skriptuni scaled\magstep0      % 

\topmargin -0.5 cm
\oddsidemargin 0cm
\evensidemargin 0cm
\textwidth 16cm
\textheight 22.2cm
\parindent 0.5cm

\input{epsf} 
\newfont{\bfsf}{cmssdc10 scaled 1000}
\newfont{\bfsfsm}{cmssdc10 scaled 800}
\newcommand{\meta}{{\sf META-Frame}}

\newcommand{\bnfterm}[1]{\underline{\tt #1}}
\newcommand{\bnfstar}{\mbox{$^*$}}
\newcommand{\bnfplus}{\mbox{$^+$}}
\newcommand{\bnfchoice}{\mbox{$\,\mid\,$}}
\newcommand{\bnfempty}{\mbox{$\lambda$}}

\setlength{\parindent}{0pt}

\nonfrenchspacing
\sloppy
\raggedright


\begin{document}

% \initfloatingfigs

%----------------------------------------------------------------------------
% TITLE AND STUFF
%----------------------------------------------------------------------------


\footskip5cm
\thispagestyle{empty}
\begin{center}
{\skriptuni UNIVERSITaTPASSAU} \\
\medskip

%\ \\
{\large {\sf Fakult\"at f\"ur Mathematik und Informatik}} \\
%{\sc D--94030 Passau} \\
\ \\
\ \\
{\fmilogofont P}
\ \\
\ \\
\ \\
\ \\
{\Large {\bf The ffgraph Library}}\\
\ \\
\ \\              
\ \\
\ \\


Carsten Friedrich\\[4mm]
Lehrstuhl f\"{u}r Informatik \\
Universit\"at Passau\\
Innstr. 33\\
D-94032 Passau (Germany)\\
tel: +49 851 509.3090 \\
fax: +49 851 509.3092 \\
{\tt friedric@fmi.uni-passau.de}\\
\ \\
\ \\
\ \\
\ \\
\ \\
\ \\
\ \\
MIP-9520 \\
Dezember 1995 \\
\end{center}
\ \\
\ \\
\ \\
\ \\
\ \\ 
\ \\ 

\newpage 
\thispagestyle{empty}

\newpage

\pagestyle{headings}

\tableofcontents

\newpage

\section{Introduction} 
\label{introduction}

\input ffIntro

\section{Installing the library} 
\label{installing}

\input ffInstall

\section{Using the graph library}  
\label{using_the_graph_library}

\input ffUsing

\section{Compiling your program}
\label{compiling_your_program}

\input ffCompile

\section{Labels}  
\label{labels}

\input ffLabels

\section{Inspectors}
\label{inspectors}

\input ffInspector

\section{Graphic contexts}
\label{graphic}

\input ffGraphic

\section{Layouters}  
\label{layouters}

\input ffLayout

\section{The important classes in detail} 
\label{classes_in_detail}

\subsection{Standard classes}
\label{standard_classes}

\input ffStdClasses

\subsection{Label classes}
\label{label_classes}

\input ffLabelClasses

\subsection{Graphic classes}
\label{graphic_classes}

\input ffGraphicClasses

\section{The ffgraph graph editor} 

A graph editor called {\bf ffgraph} is included in the
distribution to give you an example of how to use the library.
You can pass a list of graphs to {\bf ffgraph} which will be
loaded and displayed.

\begin{appendix}

% \section{Copyright information} 

% \input COPYRIGHT

\section{The extended sgraph file format}
\label{sgraph_file_format}

\input ffSgraph

\section{The class hierarchy}
\label{the_class_hierachy}

This figure shows the class hierachy of the most important
classes in the library.

\begin{figure}[htb]
\fbox{\epsfig{figure=hierachy.eps}}
\caption{The class hierarchy}
\end{figure}

\clearpage
\newpage

\section*{Acknowledgements}

I am grateful to Prof. Bernhard Steffen, Andreas Cla\ss en, Dr. Ruth
Wishart, Karlheinz Agsteiner and all of the \meta{} team for their
constructive cooperation in the design and implementation of the ffgraph
library.

\newpage 

\bibliographystyle{alpha}
\bibliography{ffgraph}

\end{appendix}

\newpage

\begin{center}

\bf{Technische Berichte\\
    der Fakult\"at f\"ur Mathematik und Informatik\\
    Universit\"at Passau ab 1993}
\end{center}

\begin{longtable}{lp{10cm}l}

\typeout{}\noindent
\typeout{   >>> MIPs JAHRGANG 93... <<<}\noindent

MIP-9301&
B.S. Chlebus, M. Kaufmann, J.F. Sibeyn: Deterministic Permutation Routing on Meshes.\\

MIP-9302&
M. M"uller: Bondgraph Execution as a New Algorithm for Circuit Simulation.\\

MIP-9303&
C. Lengauer: Loop Parallelization in the Polytope Model.\\

MIP-9304&
M. Griebl, C. Lengauer: On the Space-Time Mapping of WHILE-Loops.\\

MIP-9305&
V. Weispfenning: A New Approach to Quantifier Elimination for Real Algebra.\\

MIP-9306&
A. Kemper, G. Moerkotte: Physical Object Management.\\

MIP-9307&
M. Steinbrunn, G. Moerkotte, A. Kemper: Optimizing Join Orders.\\

MIP-9308&
A. Kemper, G. Moerkotte, K. Peithner, M. Steinbrunn: Optimizing Disjunctive Queries in Object Bases.\\

MIP-9309&
A. Kemper, D. Kossmann: Dual-Buffering Strategies in Object Bases.\\

MIP-9310&
J. Knoop, O. R"uthing, B. Steffen: Optimal Code Motion: Theory and Practice.\\

MIP-9311&
B. Freitag, T. Margaria, B. Steffen: A Pragmatic Approach to Software Synthesis.\\

MIP-9312&
O. Burkart, B. Steffen: Pushdown Processes: Parallel Composition and Model Checking.\\

MIP-9313&
B. Steffen, T. Margaria, B. Freitag: Module Configuration by Minimal Model Construction.\\

MIP-9314&
A. Geser: A Solution to Zantema's Problem.\\

MIP-9315&
S. Gorlatch, C. Lengauer: Parallelization of Divide-and-Conquer in the Bird-Meertens Formalism.\\

\typeout{}\noindent
\typeout{   >>> MIPs JAHRGANG 94... <<<}\noindent

MIP-9401&
U. Fraus: Inductive Theorem Proving for Algebraic Specifications � TIP System User's Manual.\\

MIP-9402&
O. Burkart, B. Steffen: Composition, Decomposition and Model Checking of Pushdown Processes.\\

MIP-9403&
A. Geser, J. Knoop, G. L"uttgen, O. R"uthing, B. Steffen: Chaotic Fixed Point Iterations.\\

MIP-9404&
V. Weispfenning: Parametric Linear and Quadratic Optimization by Elimination.\\

MIP-9405&
W. Bachl, F.J. Brandenburg, T. Hickl: Hierarchical Graph Design Using HiGraD.\\

MIP-9406&
M. Griebl, C. Lengauer: A Communication Scheme for the Distributed Execution of while Loops.\\

MIP-9407&
A. Geser: An Improved General Path Order.\\

MIP-9408&
S. Graf, H. Luschgy: The quantization of the Cantor measure.\\

MIP-9409&
J. Knoop, B. Steffen, J. Vollmer: Parallelism for Free: Efficient and Optimal Bitvector Analyses for Parallel Programs.\\

MIP-9410&
U. Lechner, C. Lengauer, M. Wirsing: An Object-Oriented Airport.\\

MIP-9411&
T. Margaria, M. Griva, R. Tesio: Semantic Extraction for the Automatic Verification of VHDL Descriptions.\\

MIP-9412&
M. Steinbrunn, K. Peithner, G. Moerkotte, A. Kemper: Bypassing Joins in Disjunctive Queries.\\

MIP-9413&
J. Knoop, O. R"uthing, B. Steffen: A Tool Kit for Constructing Optimal Interprocedural Data Flow Analyses.\\

MIP-9414&
C. Lengauer, M. Griebl: On the Parallelization of Loop Nests Containing while Loops.\\

MIP-9415&
G. Ritter: Jackson's Theorems and the Number of Hidden Units in Neural Networks for Uniform Approximation.\\

\typeout{}\noindent
\typeout{   >>> MIPs JAHRGANG 95... <<<}\noindent

MIP-9501&
B. Steffen: Finite Model Checking and Beyond.\\

MIP-9502&
K.G. Larsen, B. Steffen, C. Weise: A Constraint Oriented Proof Methodology Based on Modal Transition Systems.\\

MIP-9503&
V. Weispfenning, R. Xue: Parametric Mixed Integer Programming by Elimination.\\

MIP-9504&
V. Weispfenning: Solving Parametric Polynomial Equations and Inequalities by Symbolic Algorithms.\\

MIP-9505 &
S. Graf, B. Steffen, G. L"uttgen: Compositional Minimization of Finite State Systems Using Interface Specifications.\\

MIP-9506&
F.J. Brandenburg, M. Himsolt, H.J. R"oder, K. Skodinis: Designing Data Flow Graphs by Graph Grammars.\\

MIP-9507&
A. Eickler, C.A. Gerlhof, D. Kossmann: A Performance Evaluation of OID Mapping Techniques.\\

MIP-9508&
S. Gorlatch: From Transformations to Methodology in Parallel Program Development: A Case Study.\\

MIP-9509&
B. Steffen, T. Margaria, A. Cla\ss en: Heterogeneous Analysis and Verification for Distributed Systems.\\

MIP-9510&
B. Freitag: Method Resolution and Virtual Classes in a Deductive Object-Oriented Database.\\

MIP-9511&
J. Knoop, B. Steffen, J. Vollmer: ``Optimal Code Motion for Parallel Programs''.\\


MIP-9512&
S. Gorlatch: Constructing List Homomorphisms.\\

MIP-9513&
K. Donner: Ersatzarithmetiken f"ur die fehlerstabile Implementation numerischer Algorithmen am Beispiel der Polynominterpolation.\\

MIP-9514&
K. Donner: Skalierung von Matrizen und numerische Stabilit"at der Gau\ss-Elimination\\

MIP-9515&
K. Donner: Orthogonale Polynome f"ur zeit"aquidistante Abtastung von Signalen\\

MIP-9516&
B. Steffen, T. Margaria, A. Cla\ss en, V. Braun: Incremental Formalization: a Key to Industrial Success\\

MIP-9517&
A. Dolzmann, Th. Sturm: Simplification of Quantifier-free Formulas over Ordered Fields.\\

MIP-9518&
M. Mendler: A Timing Refinement of Intuitionistic Proofs and its Application to the Timing Analysis of Combinational Circuits\\

MIP-9519&
T. Margaria (Hrsg.): Programmiersprachen und Grundlagen der Programmierung\\

\end{longtable}


\end{document}

